\section{Eliminaci\'on gaussiana y determinante}
\whenuse{Resolver sistemas lineales, calcular rango o determinante; elegir la
versi\'on real con pivoteo o modular seg\'un el dominio del problema.}

Usar para resolver sistemas lineales, calcular rango o determinante. Esta
versi\'on trabaja con reales y pivoteo parcial; para m\'odulo primo reemplazar la
divisi\'on por inverso modular.

\complexity{tiempo $O(nm\min(n,m))$ para $n$ ecuaciones y $m$ variables;
memoria $O(nm)$. Determinante cuadrado: $O(n^3)$.}

\begin{lstlisting}[style=cpp]
int gaussian_elimination(vector<vector<long double>> matrix,
                         vector<long double> &solution) {
    const long double epsilon = 1e-12L;
    int rows = (int)matrix.size();
    int columns = (int)matrix[0].size() - 1;
    vector<int> pivot_row(columns, -1);
    int row = 0;
    for (int column = 0; column < columns && row < rows; ++column) {
        int selected = row;
        for (int candidate = row; candidate < rows; ++candidate)
            if (fabsl(matrix[candidate][column]) > fabsl(matrix[selected][column]))
                selected = candidate;
        if (fabsl(matrix[selected][column]) < epsilon) continue;
        swap(matrix[selected], matrix[row]);
        pivot_row[column] = row;
        for (int candidate = 0; candidate < rows; ++candidate) if (candidate != row) {
            long double factor = matrix[candidate][column] / matrix[row][column];
            for (int index = column; index <= columns; ++index)
                matrix[candidate][index] -= factor * matrix[row][index];
        }
        ++row;
    }
    for (int candidate = 0; candidate < rows; ++candidate) {
        bool all_zero = true;
        for (int column = 0; column < columns; ++column)
            all_zero &= fabsl(matrix[candidate][column]) < epsilon;
        if (all_zero && fabsl(matrix[candidate][columns]) >= epsilon) return 0;
    }
    solution.assign(columns, 0);
    for (int column = 0; column < columns; ++column)
        if (pivot_row[column] != -1)
            solution[column] = matrix[pivot_row[column]][columns]
                             / matrix[pivot_row[column]][column];
    return row == columns ? 1 : 2; // unica / infinitas
}
// matriz cuadrada es invertible:
long double determinant(vector<vector<long double>> matrix) {
    const long double epsilon = 1e-12L;
    int size = (int)matrix.size();
    long double answer = 1;
    for (int column = 0; column < size; ++column) {
        int pivot = column;
        for (int row = column; row < size; ++row)
            if (fabsl(matrix[row][column]) > fabsl(matrix[pivot][column])) pivot = row;
        if (fabsl(matrix[pivot][column]) < epsilon) return 0;
        if (pivot != column) swap(matrix[pivot], matrix[column]), answer = -answer;
        answer *= matrix[column][column];
        for (int row = column + 1; row < size; ++row) {
            long double factor = matrix[row][column] / matrix[column][column];
            for (int index = column + 1; index < size; ++index)
                matrix[row][index] -= factor * matrix[column][index];
        }
    }
    return answer;
}
\end{lstlisting}
\begin{lstlisting}[style=cpp]
vector<vector<long double>> augmented={
    {1,1,3},
    {2,-1,0}
};
vector<long double> solution;
int type=gaussian_elimination(augmented,solution);

if(type==0) cout<<"Sin solucion\n";
else if(type==1)
    cout<<solution[0]<<" "<<solution[1]<<'\n'; // 1 2
else
    cout<<"Infinitas soluciones\n";
vector<vector<long double>> coefficients={
    {1,1},
    {2,-1}
};
cout<<(double)determinant(coefficients)<<'\n'; // -3
\end{lstlisting}
